% Options for packages loaded elsewhere
% Options for packages loaded elsewhere
\PassOptionsToPackage{unicode}{hyperref}
\PassOptionsToPackage{hyphens}{url}
\PassOptionsToPackage{dvipsnames,svgnames,x11names}{xcolor}
%
\documentclass[
  letterpaper,
  DIV=11,
  numbers=noendperiod]{scrartcl}
\usepackage{xcolor}
\usepackage{amsmath,amssymb}
\setcounter{secnumdepth}{-\maxdimen} % remove section numbering
\usepackage{iftex}
\ifPDFTeX
  \usepackage[T1]{fontenc}
  \usepackage[utf8]{inputenc}
  \usepackage{textcomp} % provide euro and other symbols
\else % if luatex or xetex
  \usepackage{unicode-math} % this also loads fontspec
  \defaultfontfeatures{Scale=MatchLowercase}
  \defaultfontfeatures[\rmfamily]{Ligatures=TeX,Scale=1}
\fi
\usepackage{lmodern}
\ifPDFTeX\else
  % xetex/luatex font selection
\fi
% Use upquote if available, for straight quotes in verbatim environments
\IfFileExists{upquote.sty}{\usepackage{upquote}}{}
\IfFileExists{microtype.sty}{% use microtype if available
  \usepackage[]{microtype}
  \UseMicrotypeSet[protrusion]{basicmath} % disable protrusion for tt fonts
}{}
\makeatletter
\@ifundefined{KOMAClassName}{% if non-KOMA class
  \IfFileExists{parskip.sty}{%
    \usepackage{parskip}
  }{% else
    \setlength{\parindent}{0pt}
    \setlength{\parskip}{6pt plus 2pt minus 1pt}}
}{% if KOMA class
  \KOMAoptions{parskip=half}}
\makeatother
% Make \paragraph and \subparagraph free-standing
\makeatletter
\ifx\paragraph\undefined\else
  \let\oldparagraph\paragraph
  \renewcommand{\paragraph}{
    \@ifstar
      \xxxParagraphStar
      \xxxParagraphNoStar
  }
  \newcommand{\xxxParagraphStar}[1]{\oldparagraph*{#1}\mbox{}}
  \newcommand{\xxxParagraphNoStar}[1]{\oldparagraph{#1}\mbox{}}
\fi
\ifx\subparagraph\undefined\else
  \let\oldsubparagraph\subparagraph
  \renewcommand{\subparagraph}{
    \@ifstar
      \xxxSubParagraphStar
      \xxxSubParagraphNoStar
  }
  \newcommand{\xxxSubParagraphStar}[1]{\oldsubparagraph*{#1}\mbox{}}
  \newcommand{\xxxSubParagraphNoStar}[1]{\oldsubparagraph{#1}\mbox{}}
\fi
\makeatother


\usepackage{longtable,booktabs,array}
\usepackage{calc} % for calculating minipage widths
% Correct order of tables after \paragraph or \subparagraph
\usepackage{etoolbox}
\makeatletter
\patchcmd\longtable{\par}{\if@noskipsec\mbox{}\fi\par}{}{}
\makeatother
% Allow footnotes in longtable head/foot
\IfFileExists{footnotehyper.sty}{\usepackage{footnotehyper}}{\usepackage{footnote}}
\makesavenoteenv{longtable}
\usepackage{graphicx}
\makeatletter
\newsavebox\pandoc@box
\newcommand*\pandocbounded[1]{% scales image to fit in text height/width
  \sbox\pandoc@box{#1}%
  \Gscale@div\@tempa{\textheight}{\dimexpr\ht\pandoc@box+\dp\pandoc@box\relax}%
  \Gscale@div\@tempb{\linewidth}{\wd\pandoc@box}%
  \ifdim\@tempb\p@<\@tempa\p@\let\@tempa\@tempb\fi% select the smaller of both
  \ifdim\@tempa\p@<\p@\scalebox{\@tempa}{\usebox\pandoc@box}%
  \else\usebox{\pandoc@box}%
  \fi%
}
% Set default figure placement to htbp
\def\fps@figure{htbp}
\makeatother





\setlength{\emergencystretch}{3em} % prevent overfull lines

\providecommand{\tightlist}{%
  \setlength{\itemsep}{0pt}\setlength{\parskip}{0pt}}



 


\KOMAoption{captions}{tableheading}
\makeatletter
\@ifpackageloaded{caption}{}{\usepackage{caption}}
\AtBeginDocument{%
\ifdefined\contentsname
  \renewcommand*\contentsname{Table of contents}
\else
  \newcommand\contentsname{Table of contents}
\fi
\ifdefined\listfigurename
  \renewcommand*\listfigurename{List of Figures}
\else
  \newcommand\listfigurename{List of Figures}
\fi
\ifdefined\listtablename
  \renewcommand*\listtablename{List of Tables}
\else
  \newcommand\listtablename{List of Tables}
\fi
\ifdefined\figurename
  \renewcommand*\figurename{Figure}
\else
  \newcommand\figurename{Figure}
\fi
\ifdefined\tablename
  \renewcommand*\tablename{Table}
\else
  \newcommand\tablename{Table}
\fi
}
\@ifpackageloaded{float}{}{\usepackage{float}}
\floatstyle{ruled}
\@ifundefined{c@chapter}{\newfloat{codelisting}{h}{lop}}{\newfloat{codelisting}{h}{lop}[chapter]}
\floatname{codelisting}{Listing}
\newcommand*\listoflistings{\listof{codelisting}{List of Listings}}
\makeatother
\makeatletter
\makeatother
\makeatletter
\@ifpackageloaded{caption}{}{\usepackage{caption}}
\@ifpackageloaded{subcaption}{}{\usepackage{subcaption}}
\makeatother
\usepackage{bookmark}
\IfFileExists{xurl.sty}{\usepackage{xurl}}{} % add URL line breaks if available
\urlstyle{same}
\hypersetup{
  pdftitle={Point Estimation Study},
  colorlinks=true,
  linkcolor={blue},
  filecolor={Maroon},
  citecolor={Blue},
  urlcolor={Blue},
  pdfcreator={LaTeX via pandoc}}


\title{Point Estimation Study}
\author{}
\date{}
\begin{document}
\maketitle


\section{What is a Point Estimate?}\label{what-is-a-point-estimate}

A \textbf{point estimate} is a single numerical value used to estimate
an unknown population parameter. In real-world situations, populations
can be extremely large (or even infinite), making it impractical to
measure the true parameter directly. Instead, we rely on sample data and
compute a statistic that serves as our best guess.

For example, if we want to estimate the population mean, we use the
\textbf{sample mean} as a point estimate.

\textbf{Point Estimation} is therefore the process of using sample data
to compute this single value.

\begin{center}\rule{0.5\linewidth}{0.5pt}\end{center}

\section{What Makes a Good Point
Estimate?}\label{what-makes-a-good-point-estimate}

A good estimator should satisfy three important properties:

\begin{itemize}
\tightlist
\item
  It should be \textbf{unbiased}
\item
  It should be \textbf{precise}
\item
  It should be \textbf{accurate}
\end{itemize}

\begin{center}\rule{0.5\linewidth}{0.5pt}\end{center}

\subsection{Defining Key Terms}\label{defining-key-terms}

\subsubsection{Precision}\label{precision}

Precision refers to how consistent an estimator is across repeated
samples. If we repeat an experiment many times and get estimates that
are very close to each other, then the estimator is precise.

Mathematically, precision is related to variance:

So: - \textbf{Low variance → high precision} - \textbf{High variance →
low precision}

\begin{center}\rule{0.5\linewidth}{0.5pt}\end{center}

\subsubsection{Bias}\label{bias}

Bias measures how far an estimator is, on average, from the true
parameter value.

\begin{itemize}
\tightlist
\item
  If bias = 0 → estimator is \textbf{unbiased}
\item
  If bias ≠ 0 → estimator is \textbf{biased}
\end{itemize}

\begin{center}\rule{0.5\linewidth}{0.5pt}\end{center}

\subsubsection{Accuracy}\label{accuracy}

Accuracy combines both bias and precision. An estimator is accurate if
it is: - close to the true value (low bias) - not highly variable (low
variance)

A common way to measure accuracy is the \textbf{Mean Squared Error
(MSE)}:

This shows that even an unbiased estimator can be bad if its variance is
large.

\begin{center}\rule{0.5\linewidth}{0.5pt}\end{center}

\section{Methods of Finding Point
Estimates}\label{methods-of-finding-point-estimates}

There are two main methods we focus on:

\subsection{1. Method of Moments (MOM)}\label{method-of-moments-mom}

The Method of Moments works by matching: - \textbf{Sample moments} (from
data) - \textbf{Population moments} (from theory)

For example:

We solve these equations to estimate unknown parameters.

The justification comes from the \textbf{Law of Large Numbers}, which
tells us sample averages converge to population values as the sample
size increases.

\begin{center}\rule{0.5\linewidth}{0.5pt}\end{center}

\subsection{2. Maximum Likelihood Estimation
(MLE)}\label{maximum-likelihood-estimation-mle}

MLE takes a different approach. Instead of matching moments, it asks:

\begin{quote}
What parameter value makes the observed data most likely?
\end{quote}

The likelihood function is:

The value of that maximizes this function is the \textbf{MLE}.

\begin{center}\rule{0.5\linewidth}{0.5pt}\end{center}

\section{Simulation Study: Comparing
Estimators}\label{simulation-study-comparing-estimators}

To show real understanding, we simulate data and compare two estimators
for a Uniform distribution:

\begin{center}\rule{0.5\linewidth}{0.5pt}\end{center}

\subsection{Simulation Setup}\label{simulation-setup}




\end{document}
